\subsection{Motivación y posición cronológica}

Checkpoint~04C.1 se ejecutó después de haber observado la superioridad de métodos
locales/grafo y después de cerrar la hipótesis SIAP. Se mantiene la etiqueta 04C porque
representa la rama de modelos globales de mayor compute prevista desde el diseño de
Checkpoint~04. La pregunta no es volver a buscar el mejor modelo global, sino medir si
CatBoost aporta un error suficientemente diferente para mejorar un predictor local mediante
un peso pequeño.

La representación se restringe a C1 agronomic, con 350 features, porque 03C.2 ya había
identificado C1+CatBoost como el anchor boosting global más relevante. Reabrir cientos de
configuraciones sobre los mismos 138 labels habría aumentado el riesgo de overfitting de
selección sin aportar información independiente.

\subsection{Candidate set congelado}

Se reutilizan tres configuraciones presentes en el grid histórico de 03C.2:
\begin{table}[H]
\centering
\caption{Candidatos CatBoost de 04C.1.}
\label{tab:cp04c-cat}
\begin{tabular}{rrrr}
\toprule
depth & learning rate & $L_2$ leaf reg & iterations \\
\midrule
4 & 0.03 & 3 & 400 \\
4 & 0.06 & 6 & 400 \\
5 & 0.06 & 3 & 400 \\
\bottomrule
\end{tabular}
\end{table}

Cada configuración se entrena con seeds 42, 314 y 2718; la predicción usada para evaluación
es
\[
\hat y_q^{CB}
=
\frac{1}{3}
\sum_{s\in\{42,314,2718\}}
\hat y_{q,s}^{CB}.
\]
La ejecución se hace en GPU; no existe fallback automático a CPU. Esta decisión es de
reproducibilidad operacional, no estadística.

\subsection{Blends predefinidos}

Para el candidato estable depth 4, learning rate 0.03, $L_2=3$, se definen pesos pequeños
\[
\hat y^{(w)}
=
(1-w)\hat y^{Local}
+w\hat y^{CB},
\qquad
w\in\{0.10,0.20,0.30\}.
\]
También se construyen blends equivalentes con Graph04D como baseline. Los pesos son
deliberadamente discretos: dado que CatBoost standalone ya es más débil, la hipótesis plausible
es corrección/diversidad, no dominancia.

\subsection{Resultados target-matched}

\begin{table}[H]
\centering
\caption{Resultados principales de 04C.1.}
\label{tab:cp04c-results}
\small
\begin{tabular}{lrrr}
\toprule
Método & RMSE medio & pooled RMSE & pooled MAE \\
\midrule
Local + CatBoost 10\% & 0.487842 & 0.495915 & 0.342266 \\
Local04D & \textbf{0.487845} & \textbf{0.495741} & 0.342501 \\
Local + CatBoost 20\% & 0.488501 & 0.496728 & 0.342952 \\
Local + CatBoost 30\% & 0.489822 & 0.498177 & 0.344425 \\
Graph + CatBoost 30\% & 0.491439 & 0.500394 & 0.339809 \\
Graph04D & 0.494881 & 0.503346 & 0.340288 \\
CatBoost d5/lr.06/l2=3 & 0.515282 & 0.523869 & 0.370873 \\
CatBoost d4/lr.03/l2=3 & 0.516796 & 0.525414 & 0.371216 \\
CatBoost d4/lr.06/l2=6 & 0.520583 & 0.528819 & 0.373011 \\
\bottomrule
\end{tabular}
\end{table}

El blend 10\% parece superar Local04D en la sexta cifra decimal:
\[
0.4878453167-0.4878418585
\approx3.46\times10^{-6}.
\]
Pero esta diferencia carece de evidencia práctica. El pooled RMSE es peor:
0.495915 frente a 0.495741. Seleccionar el blend sólo por el mínimo de la media de los mismos
16 splits sería un ejemplo clásico de sobreinterpretación de ruido de validación.

\subsection{Gate leave-one-split-out}

Se restringe el universo de selección a:
\[
\{\text{Local},\ \text{CatBoost},\
\text{Local+CB10},\text{Local+CB20},\text{Local+CB30}\}.
\]
Para cada holdout $h$, el método se selecciona por RMSE medio sobre los otros 15 splits y se
evalúa en $h$. El selector escoge Local puro en 8/16 y Local+CB10 en 8/16. El resultado
concatenado es
\[
\overline{\RMSE}_{LOSO}=0.488884,
\qquad
\RMSE_{\mathrm{pooled}}=0.496838.
\]
Ambos son peores que Local04D fijo. Por tanto, la minúscula diferencia de la tabla completa no
sobrevive selección split-excluded.

\subsection{Diversidad de residuos}

Para pares de métodos se calcula correlación de residuos sobre las 656 pseudo-predicciones
target-matched:
\[
\rho(e^{Local},e^{Graph})=0.9707,
\]
\[
\rho(e^{Local},e^{CB})=0.9423,
\qquad
\rho(e^{Graph},e^{CB})=0.9143.
\]
CatBoost es algo más diverso respecto al grafo, pero sus errores siguen altamente
correlacionados. Su pérdida de accuracy standalone es demasiado grande para que esa diversidad
compense de forma robusta.

\subsection{Stress tests}

Graph+CatBoost puede mejorar algunas métricas target/state-random, pero no conserva la mejor
robustez municipality-grouped. Por ejemplo, Graph+CB30 alcanza 0.4914 target-matched y
0.4627 state-random, frente a 0.4949 y 0.4679 del grafo; sin embargo, grouped empeora de
0.5279 a aproximadamente 0.5508. Un peso 10\% casi preserva grouped, pero sigue siendo peor
que Local en el protocolo primario.

\subsection{Decisión}

\begin{decision}
CatBoost no entra como componente obligatorio de 04F. Se conserva como diagnóstico de
diversidad y sensibilidad. La razón no es que boosting sea débil en general: la evidencia
específica del problema muestra que, después de aprovechar localidad espacial, el anchor
CatBoost no aporta una mejora reproducible.
\end{decision}
